\documentclass[12pt, twoside]{article}
\input{../preamble}

\fancyhead[LE]{\thepage}
\fancyhead[RO]{Markscheme}
\fancyhead[LO]{La Scuola d'Italia / Dr.\ Huson / IB Math \\* 28 April 2026}

\begin{document}
\subsubsection*{11.13 Classwork: Algebra (no calculator) --- Markscheme}

\begin{enumerate}[itemsep=0.3cm]

%--- Problem 1: u = ln x ---
\item[1.] [Maximum mark: 4]

Solve $2(\ln x)^2 = 7\ln x - 3$.

\medskip

Let $u = \ln x$. Then $2u^2 - 7u + 3 = 0$. \hfill \textbf{M1}

$(2u - 1)(u - 3) = 0$

$u = \dfrac{1}{2}$ or $u = 3$ \hfill \textbf{A1}

$\ln x = \dfrac{1}{2} \;\Rightarrow\; x = e^{1/2} = \sqrt{e}$ \hfill \textbf{A1}

$\ln x = 3 \;\Rightarrow\; x = e^3$ \hfill \textbf{A1}

\bigskip
\hrule
\bigskip

%--- Problem 2: u = ln x with log property ---
\item[2.] [Maximum mark: 5]

Solve $\ln(x^2) = (\ln x)^2 - 3$ for $x > 0$.

\medskip

$\ln(x^2) = 2\ln x$ \quad (for $x > 0$) \hfill \textbf{M1}

Let $u = \ln x$:

$2u = u^2 - 3$

$u^2 - 2u - 3 = 0$ \hfill \textbf{M1}

$(u - 3)(u + 1) = 0$

$u = 3$ or $u = -1$ \hfill \textbf{A1}

$\ln x = 3 \;\Rightarrow\; x = e^3$ \hfill \textbf{A1}

$\ln x = -1 \;\Rightarrow\; x = e^{-1} = \dfrac{1}{e}$ \hfill \textbf{A1}

\bigskip
\hrule
\bigskip

%--- Problem 3: u = e^x ---
\item[3.] [Maximum mark: 5]

Solve $3e^{2x} + 2 = 7e^x$.

\medskip

Let $u = e^x$, so $e^{2x} = u^2$:

$3u^2 + 2 = 7u$

$3u^2 - 7u + 2 = 0$ \hfill \textbf{M1}

$(3u - 1)(u - 2) = 0$ \hfill \textbf{M1}

$u = \dfrac{1}{3}$ or $u = 2$ \hfill \textbf{A1}

$e^x = \dfrac{1}{3} \;\Rightarrow\; x = \ln \dfrac{1}{3} = -\ln 3$ \hfill \textbf{A1}

$e^x = 2 \;\Rightarrow\; x = \ln 2$ \hfill \textbf{A1}

\newpage

%--- Problem 4: trig with double-angle ---
\item[4.] [Maximum mark: 6]

Solve $\cos 2x = 3\cos x - 2$ for $0 \le x \le 2\pi$.

\medskip

Use $\cos 2x = 2\cos^2 x - 1$:

$2\cos^2 x - 1 = 3\cos x - 2$ \hfill \textbf{M1}

$2\cos^2 x - 3\cos x + 1 = 0$

Let $u = \cos x$:

$(2u - 1)(u - 1) = 0$ \hfill \textbf{M1}

$u = \dfrac{1}{2}$ or $u = 1$

$\cos x = \dfrac{1}{2}$: \quad $x = \dfrac{\pi}{3}$ \hfill \textbf{A1}

\hphantom{$\cos x = \dfrac{1}{2}$:\quad} $x = \dfrac{5\pi}{3}$ \hfill \textbf{A1}

$\cos x = 1$: \quad $x = 0$ \hfill \textbf{A1}

\hphantom{$\cos x = 1$:\quad} $x = 2\pi$ \hfill \textbf{A1}

\bigskip
\hrule
\bigskip

%--- Problem 5: u = 2^x with index manipulation ---
\item[5.] [Maximum mark: 5]

Solve $4^x + 16 = 5 \cdot 2^{x+1}$.

\medskip

$4^x = (2^x)^2$ and $5 \cdot 2^{x+1} = 10 \cdot 2^x$.

Let $u = 2^x$:

$u^2 + 16 = 10u$

$u^2 - 10u + 16 = 0$ \hfill \textbf{M1}

$u = \dfrac{10 \pm \sqrt{100 - 64}}{2} = \dfrac{10 \pm 6}{2}$ \hfill \textbf{M1}

$u = 2$ or $u = 8$ \hfill \textbf{A1}

$2^x = 2 \;\Rightarrow\; x = 1$ \hfill \textbf{A1}

$2^x = 8 \;\Rightarrow\; x = 3$ \hfill \textbf{A1}

\emph{Accept factoring $(u - 2)(u - 8) = 0$ in place of the quadratic formula for the second M1.}

\bigskip
\hrule
\bigskip

%--- Problem 6: u = sqrt(x) ---
\item[6.] [Maximum mark: 4]

Solve $2x + 3 = 7\sqrt{x}$.

\medskip

Let $u = \sqrt{x}$ (so $u \ge 0$ and $x = u^2$):

$2u^2 + 3 = 7u$

$2u^2 - 7u + 3 = 0$ \hfill \textbf{M1}

$(2u - 1)(u - 3) = 0$ \hfill \textbf{M1}

$u = \dfrac{1}{2}$ or $u = 3$

$\sqrt{x} = \dfrac{1}{2} \;\Rightarrow\; x = \dfrac{1}{4}$ \hfill \textbf{A1}

$\sqrt{x} = 3 \;\Rightarrow\; x = 9$ \hfill \textbf{A1}

\end{enumerate}
\end{document}
